\SourceLeaf{141}
\begin{quote}
Vielmehr scheint mir noch
\begin{list}{}{\setlength{\leftmargin}{1.8em}\setlength{\labelwidth}{1.4em}\setlength{\labelsep}{.4em}}
\item[3)] möglich und dienlich, dass es jetzt zur beständigen Norm gemacht werden dürfe, die Witwenpensionen jedesmal um 10 Thl. zu erhöhen, so oft sich der Fundus um 5000 Thl. vermehrt hat und die Zahl der Witwen noch nicht über das maximum von 15 gestiegen ist. Es ist klar, dass man dies thun kann, denn eine Erhöhung von 10 Thl. für 15 Witwen beträgt 150 Thl. und 5000 Thl. zu 3 proc. geben eben so viel Interesse. Dass aber der nächste Erhöhungs-Termin bald eintreten kann, wenn auch unsere Kasse keine ausserordentliche Zuflüsse erhält, dies lässt sich wenigstens sehr wahrscheinlich berechnen. Da wir gegenwärtig nur 10 Witwen zu pensioniren haben, so müssen, wenn der neue Zuschuss zu den Beiträgen bewilligt wird, alle Jahr über 1000 Thl. der Kasse bleiben, folglich 5000 Thl. schon in 5 Jahren zum fundus hinzugekommen sein, wenn sich die Zahl der Witwen indessen nicht vermehrt; setzt man aber auch den höchst unwahrscheinlichen Fall, dass die Zahl jedes Jahr um Eine Witwe vermehrt würde, bis sie das maximum von 15 erreicht hätte, so würde es doch kaum 10 Jahre anstehen können.
\end{list}
Gut möchte es wenigstens sein, wenn in dem Bericht an Kön. Regierung dieser Umstand erwähnt würde.’
\end{quote}

\noindent\hspace*{\parindent}Da über den wesentlichen Inhalt dieses Artikels weiter unten bei der Prüfung des \textsc{Kritter}’schen Gutachtens und der daran geknüpften Folgerungen, und an andern Stellen das Nöthige vorkommen wird, so sollen hier nur ein paar Nebenumstände berührt werden.

I. Auffallend ist, aus der Feder des Curators der Witwenkasse, die \textit{unrichtige} Angabe der Witwen-, oder vielmehr Pensionenzahl. Es waren damals nicht 10 Witwen, und auch nicht 10 Pensionen, sondern 8 Witwen, und, unter Hinzuzählung der Kinder des am 8. Junius 1794 verstorbenen Professors B., zusammen 9 Pensionen. Auch lässt sich diese Unrichtigkeit nicht etwa dadurch erklären, dass der Abgang der zuletzt an fremdem Orte (Halle) verstorbenen Witwe (M.) dem Curator damals noch unbekannt gewesen sei; denn es findet sich, dass der Betrag der Pension für das letzte halbe Jahr (Michaelis 1793 bis Ostern 1794) auf eine von P. mitunterzeichnete und vom 3. Mai 1794 datirte Quittung der Erbin erhoben ist.

II. Die Schlusszeile (‘Gut möchte es wenigstens u. s. w.’) lässt uns etwas im Dunkeln rücksichtlich der Frage, für was der abgeschriebene Artikel eigentlich genommen werden soll, ob für einen formlichen zur Beschlussnahme verstellten Antrag, oder nur für eine hingeworfene Idee. Von einer für alle künftigen Zeiten geltenden bestimmten Normirung der veränderlichen Pensionshöhe war weder in dem Rescript, noch in den vorhergehenden Verhandlungen die Rede gewesen. Es war dies also ein \textit{vierter} zu den bereits in Deliberation begriffenen neu hinzukommender Gegenstand, und zwar ein solcher, der den drei andern an Wichtigkeit keinesweges nachstehend die sorgfältigste allseitige Prüfung erforderte. Wer einen auf ein solches Ziel eigens gerichteten Antrag einbringt, ist sich doch der Wichtigkeit der Sache bewusst, welche durch die Bezeichnung einer solchen Lebensfrage mit ‘dieser Umstand’ schwerlich genug hervortritt. Bei einer Äusserung hingegen, die nur den Charakter eines gelegentlich hingeworfenen Gedankens hat, ist man schon nachsichtiger gegen eine durch Unachtsamkeit entschlüpfte Unrichtigkeit, und gegen eine noch mangelhafte Wortfassung. Durch die Numerirung mit (3) lasse man sich hiebei nicht irre machen. P. \textit{zählt} in seinem Aufsatze nicht die Vorschläge, sondern die aus den vorausgeschickten Rechnungsüberschlägen von ihm abgeleiteten Folgerungen. Sein Nr. 1 enthält \textit{die} Folgerung, dass alle gegen die Verlängerung der Waisenpensionen vorgebrachten oder vorzubringenden Einwürfe sich erledigen, wenn man die proponirte Erhöhung der Beiträge annehme; und sein Nr. 2 die, dass zwar die allgemeine Erhöhung der Pensionen auf 130 Rth. füglich sofort geschehen könne, die exceptionelle Erhöhung auf 150 Rth. für die sechs ältesten Witwen hingegen weder gerecht noch rathsam sei.
